\chapter*{不定积分每日一题}

\refer{《积分的方法与技巧》金玉明\ 顾新身\ 毛瑞庭\ /\ 编著}

\exercise{求积分$\int{\frac{x}{\sqrt{1+\sqrt[3]{x^2}}}\d x}$}%在花括号内输入题目

\teacher{张晓鹏}%在花括号内输入你的名字

\postdate{2023年4月5日}%在花括号内输入你想要发布那一天的日期

\result{\frac{3}{5}\qty(\sqrt{1+\sqrt[3]{x^2}})^5-2\qty(\sqrt{1+\sqrt[3]{x^2}})^3+3\sqrt{1+\sqrt[3]{x^2}}+C}%在花括号内输入本题结果

\analysis{
    对于这样的无理代数函数的积分，我们往往通过换元法转化到，有理分式积分或者三角积分（三角积分）去解决.下面是这题的解法：

    令$y=\sqrt{1+\sqrt[3]{x^2}}$，则我们有
    \formula{1+\sqrt[3]{x^2}=y^2,\ x=\qty(y^2-1)^{\frac{3}{2}},\ \d x=3y\qty(y^2-1)^{\frac{1}{2}}\d y}
    代入原积分中，我们有
    \formula{
        \begin{aligned}
            \int{\frac{x}{\sqrt{1+\sqrt[3]{x^2}}}\d x}&=\int{\frac{\qty(y^2-1)^{\frac{3}{2}}\cdot 3y\qty(y^2-1)^{\frac{1}{2}}\d y}{y}}=3\int{\qty(y^2-1)^2\d y}\\
            &=3\int{\qty(y^4-2y^2+1)\d y}=\frac{3}{5}y^5-2y^3+3y+C\\
            &=\frac{3}{5}\qty(\sqrt{1+\sqrt[3]{x^2}})^5-2\qty(\sqrt{1+\sqrt[3]{x^2}})^3+3\sqrt{1+\sqrt[3]{x^2}}+C
        \end{aligned}
    }
    换元法结束，务必记得代回去嗷.
}%在花括号内输入分析